% !PRINTED 276
\weberpage{276}
%<<CONT
gleichfalls ein Zeichenwechsel verloren gehen; daraus schliessen wir auf den Satz:

\gesperrt{Zwischen zwei auf einander folgenden Wurzeln von $P_m=0$ liegt eine und nur eine Wurzel von $P_{m-1}=0$.}

\begin{center}
\S.\,88.

\gesperrt{Zweites Beispiel.}
\end{center}

Als zweites Beispiel betrachten wir eine Gleichung, die wir der Kürze wegen die \gesperrt{Säcular-Gleichung} nennen wollen, weil die Untersuchung der säcularen Störungen der Planeten zuerst auf sie geführt hat\footnote{\gesperrt{Laplace}, Histoire de l'Acad\'emie des Sciences 1772. Von neueren Werken kann man dar\"uber vergleichen \gesperrt{Dziobek}, Theorie der Planetenbewegung, Leipzig 1888.}). Sie ist auch sonst wohl unter diesem Namen bekannt, kommt aber auch in vielen anderen Untersuchungen vor, z. B. bei der Bestimmung der Hauptaxen einer Fläche zweiten Grades, in der Theorie der kleinen Schwingungen; ihre allgemeine analytische Bedeutung liegt darin, dass sie eine besondere, die sogenannte orthogonale Transformation einer quadratischen Form in eine Summe von Quadraten liefert. Wir wollen hier die Gleichung nehmen, wie sie vorliegt, ohne Beziehung auf irgend eine Anwendung, und wollen sie nach unseren allgemeinen Sätzen discutiren.

Es sei $a_{i,k}$, wenn $i$ und $k$ die Reihe der Indices $1, 2, 3 \ldots n$ durchlaufen, irgend ein System reeller Grössen und

\weq{1}{a_{i,k} = a_{k,i}}

vorausgesetzt. Wir betrachten die symmetrische Determinante

\weq{2}{L_n(x) = \begin{vmatrix} a_{1,1}-x, & a_{1,2} & \ldots & a_{1,n} \\ a_{2,1}, & a_{2,2}-x & \ldots & a_{2,n} \\ \cdot & \cdot & \cdot & \cdot \\ a_{n,1}, & a_{n,2} & \ldots & a_{n,n}-x \end{vmatrix}}

die eine ganze rationale \gesperrt{Function} $n^{\text{ten}}$ Grades von $x$ ist, in der $x^n$ den Co\"efficienten $(-1)^n$ hat. Die Gleichung $L_n(x)=0$ soll der Gegenstand unserer Betrachtung sein.

Wir bilden die mit abwechselnden Zeichen genommene Kette der Haupt-Unterdeterminanten

\weq{3}{L_n(x), \; -L_{n-1}(x), \; L_{n-2}(x) \ldots, \; (-1)^{n-1}L_1(x), \; (-1)^n,}
%CONT>>
